\section{Solutions for Chapter~\ref{ch:centers}}

\paragraph{Exercise~\ref{ex:Z3-halfbraid}.}
\phantomsection
\label{sol:ex:Z3-halfbraid}
The scalar $\sigma_1$ must obey $\sigma_1^3=1$, and
$\sigma_h=\sigma_1^h$.  Its three choices are $\omega^k$,
giving the labels $(g,k)$ for all nine pairs.  The twist is
$\omega^{gk}$.  Five pairs have $gk=0$: the three with $g=0$
and the two further pairs with $k=0$.  Two pairs have product one,
and two have product two.  Hence the twist multiplicities are
$5,2,2$ for $1,\omega,\omega^2$.



\paragraph{Exercise~\ref{ex:Z3-tube}.}
\phantomsection
\label{sol:ex:Z3-tube}
If $g\ne g'$, multiplication vanishes.  Otherwise set $t=h+h'$
in the double sum:
\begin{align*}
 p_{g,k}p_{g,l}
 &=\frac19\sum_{t,h'}\omega^{-k(t-h')-lh'}e_{g,t}\\
 &=\frac19\sum_t\omega^{-kt}e_{g,t}
             \sum_{h'}\omega^{(k-l)h'}
 =\delta_{kl}p_{g,k}.
\end{align*}
Summing over $k$ kills every $h\ne0$, so
$\sum_kp_{g,k}=e_{g,0}$.  Summing over $g$ gives the algebra
unit.  Only $p_{0,k}$ occur in the untwisted sector.



\paragraph{Exercise~\ref{ex:tube171}.}
\phantomsection
\label{sol:ex:tube171}
For a fixed invertible $c$, tensoring on either side permutes the
six simples.  For each $a$ there is exactly one $b$ giving a
one-dimensional morphism space, so these choices contribute
$3\cdot6=18$ in total.

Fix now $c=R_l$. If $a=A_i,b=A_j$, then
$ca=R_{l-i}$ and $bc=R_{j+l}$, so a morphism exists exactly
when $j=-i$. There are three such one-dimensional spaces.
If exactly one of $a,b$ is
noninvertible, each of the nine pairs contributes one, giving
nine in each orientation.  If both are noninvertible, the two
products share all three noninvertible summands, contributing
$9\cdot3=27$. For $a=R_i,b=R_j$, their invertible summands
are $A_{l-i}$ and $A_{j-l}$; these agree exactly when $j=2l-i$.
There are three such pairs, adding three. Thus each noninvertible $c$
contributes $3+9+9+27+3=51$, and
\[
 \dim\operatorname{Tube}(\HH_3)=18+3\cdot51=171.
\]
This independently checks the matrix-block sum in
\eqref{eq:H3-tube-Wedderburn}.



\paragraph{Exercise~\ref{ex:forgetful-qdim}.}
\phantomsection
\label{sol:ex:forgetful-qdim}
The forgetful functor preserves the positive categorical dimension.
The images have dimensions
$1$, $1+3d=d^2$, $2+3d=s$, and $3d$, respectively.
For the center object $a$, its ordinary tube module has dimension
$2+1+1+1=5$, while its categorical dimension is $2+3d$.
The former counts summand multiplicities; the latter weights each
noninvertible summand by its categorical dimension $d$.  They are
different trace constructions and need not agree.



\paragraph{Exercise~\ref{ex:all-local}.}
\phantomsection
\label{sol:ex:all-local}
For a right action $r_M:M\otimes A\to M$, locality is the extra
equation
\[
 r_M\circ c_{A,M}\circ c_{M,A}=r_M.
\]
An ordinary module satisfies the associative action equations;
it need not satisfy this double-braiding equation. If $A$ is
Lagrangian in a positive modular category $\mathcal C$, as here, then
\[
 \operatorname{Dim}(\mathcal C_A^0)
   =\frac{\operatorname{Dim}(\mathcal C)}{\qdim(A)^2}=1,
 \qquad
 \operatorname{Dim}(\mathcal C_A)
   =\frac{\operatorname{Dim}(\mathcal C)}{\qdim(A)}=\qdim(A).
\]
Thus the local category is $\mathrm{Vec}$, while the all-module
category can be nontrivial. For the canonical algebra in
$Z(\HH_3)$ these categories are respectively $\mathrm{Vec}$ and
$\HH_3$. Holomorphicity concerns the first category.
The decomposition $A=\bigoplus_i m_iX_i$ gives an object and its
multiplicities. Multiplication still requires actual maps
$\mu:A\otimes A\to A$ and a unit, satisfying associativity and
commutativity with the specified associator and braiding. Those
maps are not contained in the list of summands.

